\documentclass{tietreport}

% ============================================================
% PACKAGES
% ============================================================

\usepackage{graphicx}
\usepackage{float}
\usepackage{booktabs}
\usepackage{tabularx}
\usepackage{array}
\usepackage{enumitem}
\usepackage{hyperref}

% Fix fancyhdr head-height warning from the report template
\setlength{\headheight}{15pt}

% ============================================================
% PROJECT INFORMATION
% ============================================================

% Course / Lab header
\TitlePageHeader{%
    UCS503P - Software Engineering%
}

% Main project title
\ProjectTitle{%
    AI-Enabled Smart Library Management\\
    and Resource Recommendation System%
}

% Report type
\ProjectSubTitle{%
    Prototype Report%
}

% Student information
\author{%
    \texttt{1024160010} Arnav Agarwal \\
    \texttt{1024160135} Aradhya Goyal \\
    \texttt{1024160025} Amanjot Kaur Sidhu%
}

% Degree / programme
\TitlePageSubText{%
    BE Third Year, CSE%
}

% Faculty advisor
\AdvisorName{%
    Dr. Jeelani Asif%
}

% Evaluation information
\TitlePageFooterText{{%
    \large Prototype Evaluation \\[0.2cm]
    \bfseries Computer Science and Engineering
    Department%
    } \\[0.15cm]
    Thapar Institute of Engineering and Technology,
    Patiala%
}

% ============================================================
% DOCUMENT
% ============================================================

\begin{document}

% ============================================================
% TITLE PAGE
% ============================================================

\maketitle

% ============================================================
% TABLE OF CONTENTS
% ============================================================

\tableofcontents

% ============================================================
% LIST OF FIGURES
% ============================================================

\listoffigures

% ============================================================
% CHAPTER 1: INTRODUCTION
% ============================================================

\chapter{Introduction}

\section{Background and Context}

University libraries manage large collections of books and other learning
resources while handling student records, issue and return transactions,
overdue items, fines, and resource availability. Conventional library
management systems primarily provide storage, search, and CRUD-based
operations, leaving several administrative activities manual and reactive.

The proposed Smart Library Management System combines conventional library
workflows with a lightweight multi-agent artificial intelligence layer. The
system is designed to assist students with resource discovery and
recommendation while helping administrators with record onboarding,
overdue monitoring, notifications, and resource-demand analysis.

The project uses artificial intelligence as an enabling technology rather
than as a research objective. The primary focus remains dependable
software engineering, workflow automation, system integration, measurable
improvement, and practical deployment.

\section{Problem Statement}

University libraries handle a large number of resource records, student
records, borrowing transactions, and overdue items. Although conventional
library systems can store and retrieve these records, several operational
problems remain.

\begin{itemize}

    \item \textbf{Manual resource discovery:}
    Students may need to navigate through multiple search and filtering
    options to determine whether a suitable resource is available.

    \item \textbf{Manual record onboarding:}
    Importing large numbers of student or member records can require
    repetitive validation and manual duplicate checking.

    \item \textbf{Reactive overdue management:}
    Overdue resources may only be identified after staff manually inspect
    issue and due-date records.

    \item \textbf{Limited demand visibility:}
    Library administrators may not have an immediate way to identify
    frequently issued resources with insufficient available copies.

    \item \textbf{Data without assistance:}
    A conventional database stores library information but does not
    directly assist users in converting that information into actionable
    answers or recommendations.

\end{itemize}

\textbf{Why this matters now:}
A university library serves a large student population with continuously
changing resource requirements. Improving resource discovery, record
processing, overdue monitoring, and demand visibility can reduce repetitive
administrative effort and improve the usability of existing library
resources.

\section{Project Objectives}

The prototype is designed to achieve the following objectives:

\begin{enumerate}

    \item Develop a functional web-based Smart Library Management System.

    \item Provide resource and student/member management.

    \item Support resource search and real-time availability display.

    \item Implement issue, return, due-date, and fine workflows.

    \item Provide natural-language resource discovery using the Library
    Intelligence Agent.

    \item Provide bulk record validation and duplicate detection using the
    Record Onboarding Agent.

    \item Provide automated overdue monitoring and notification support
    using the Library Monitoring Agent.

    \item Analyse resource demand and availability for administrative
    decision support.

    \item Maintain controlled interaction between AI agents and the
    underlying database.

    \item Support automated testing, continuous integration, and
    deployment.

\end{enumerate}

% ============================================================
% CHAPTER 2: METHODOLOGY AND DESIGN
% ============================================================

\chapter{Methodology and Design}

\section{Development Methodology}

The project follows an incremental engineering approach based on rapid
time-to-value and short feedback cycles.

The initial implementation focuses on establishing a usable minimum viable
product containing authentication, resource management, resource search,
issue and return workflows, and database integration.

AI-assisted functionality is introduced incrementally after the core
library workflow is operational. Continuous integration is established
early so that code changes can be checked through automated testing,
linting, and build validation.

Each development iteration produces a demonstrable increment that can be
tested independently and improved using feedback.

\section{System Architecture}

The system follows a modular three-tier web architecture integrated with a
controlled AI agent layer.

\subsection{High-Level Design}

The major components of the system are:

\begin{enumerate}

    \item \textbf{Student / Administrator Web Interface:}
    Provides interfaces for students and administrators to access library
    services and management functions.

    \item \textbf{Backend Application and REST APIs:}
    Handles authentication, authorization, library workflows, resource
    management, student records, transactions, and communication with the
    AI agents.

    \item \textbf{MySQL Database:}
    Stores library resources, student/member records, borrowing
    transactions, due dates, fines, and availability information.

    \item \textbf{Controlled AI Agent Layer:}
    Provides the Library Intelligence Agent, Record Onboarding Agent, and
    Library Monitoring Agent through controlled application APIs and
    tools.

\end{enumerate}

The architecture separates deterministic operations from AI-assisted
interpretation and recommendation.

Critical operations such as issue and return transactions, availability
calculation, due-date checking, and fine calculation are handled by the
backend and database layers.

\subsection{Technical Stack}

The project uses mature technologies suitable for rapid engineering and
deployment.

\begin{itemize}

    \item \textbf{Frontend:} React.js / Web Client

    \item \textbf{Backend:} REST-based application service

    \item \textbf{Database:} MySQL

    \item \textbf{AI Layer:} Controlled AI agents with structured outputs

    \item \textbf{Communication:} REST APIs and structured JSON

    \item \textbf{Testing:} Automated unit and integration testing

    \item \textbf{CI/CD:} Automated build, test, and deployment pipeline

    \item \textbf{Deployment:} Containerized or cloud-hosted application

\end{itemize}

% ============================================================
% SYSTEM DIAGRAMS
% ============================================================

\section{System Design and Diagrams}

\subsection{Use Case Diagram}

The Use Case Diagram defines the interactions between the primary human
actors, Students and Librarian, and the three specialized AI agents.

The system uses the following final agent terminology:

\begin{itemize}

    \item Library Intelligence Agent

    \item Record Onboarding Agent

    \item Library Monitoring Agent

\end{itemize}

\begin{figure}[H]
    \centering
    \includegraphics[
        width=0.82\textwidth,
        keepaspectratio
    ]{docs/Diagrams/Use Case Diagram.png}
    \caption{Use Case Diagram of the Smart Library Management System.}
    \label{fig:usecase}
\end{figure}

% ============================================================
% DATA FLOW DIAGRAMS
% ============================================================

\subsection{Data Flow Diagram}

The Data Flow Diagrams describe the movement of information between
students, librarians, the Smart Library Management System, the database,
and the controlled AI agent layer.

\subsubsection{Level 0 -- Context Diagram}

The Level 0 Context Diagram presents the overall system boundary and the
major information exchanges between users, AI agents, and the Smart Library
Management System.

\begin{figure}[H]
    \centering
    \includegraphics[
        width=0.88\textwidth,
        keepaspectratio
    ]{docs/Diagrams/DFD_Level0.png}
    \caption{Level 0 Context Data Flow Diagram.}
    \label{fig:dfd0}
\end{figure}
\subsubsection{Level 1 -- System Subsystems}

The Level 1 DFD decomposes the Smart Library Management System into its
major functional subsystems, including intelligent search, library records
and resource management, circulation and fine tracking, and monitoring and
notification.

\begin{figure}[H]
    \centering
    \includegraphics[
        width=0.88\textwidth,
        keepaspectratio
    ]{docs/Diagrams/DFD_Level1.png}
    \caption{Level 1 Data Flow Diagram showing the major system
    subsystems.}
    \label{fig:dfd1}
\end{figure}

\subsubsection{Level 2 -- Intelligent Search Process}

The Level 2 DFD provides a detailed representation of the intelligent
resource-search workflow, including natural-language query processing,
catalogue retrieval, availability checking, and recommendation generation.

\begin{figure}[H]
    \centering
    \includegraphics[
        width=0.88\textwidth,
        keepaspectratio
    ]{docs/Diagrams/Level_2_DFD.png}
    \caption{Level 2 Data Flow Diagram for Intelligent Resource Search.}
    \label{fig:dfd2}
\end{figure}

% ============================================================
% ENTITY RELATIONSHIP DIAGRAM
% ============================================================

\subsection{Entity Relationship Diagram}

The Entity Relationship Diagram represents the principal database entities
and their relationships.

The major entities include:

\begin{itemize}

    \item STUDENTS

    \item EMPLOYEE

    \item ASSETS

    \item BOOKS

    \item ISSUING\_DETAILS

    \item FINES

    \item NOTIFICATION\_QUEUE

    \item SEARCH\_DEMAND\_LOGS

\end{itemize}

\begin{figure}[H]
    \centering
    \includegraphics[
        width=0.92\textwidth,
        keepaspectratio
    ]{docs/Diagrams/ER_Diagram.png}
    \caption{Entity Relationship Diagram of the Smart Library Management
    System.}
    \label{fig:er}
\end{figure}

% ============================================================
% ACTIVITY / SWIMLANE DIAGRAM
% ============================================================

\subsection{Activity and Swimlane Diagram}

The Activity and Swimlane Diagram illustrates the interaction between the
student, web interface, backend application, Library Intelligence Agent,
MySQL database, and librarian/notification layer during intelligent
resource discovery and borrowing.

\begin{figure}[H]
    \centering
    \includegraphics[
        width=0.96\textwidth,
        keepaspectratio
    ]{docs/Diagrams/SwimLane_Diagram.png}
    \caption{Activity and Swimlane Diagram for Intelligent Resource
    Discovery and Borrowing.}
    \label{fig:swimlane}
\end{figure}

% ============================================================
% DATABASE DESIGN
% ============================================================

\section{Database Design}

The system uses MySQL as the primary relational database.

The existing core entities include student records, employee records,
library assets, books, issuing transactions, and fines.

The database is extended with two additional entities:

\begin{itemize}

    \item \textbf{NOTIFICATION\_QUEUE:}
    Stores notifications generated for overdue and other
    library-management events.

    \item \textbf{SEARCH\_DEMAND\_LOGS:}
    Stores natural-language search queries and extracted categories for
    subsequent resource-demand analysis.

\end{itemize}

The database design uses primary keys, foreign-key relationships, and
indexes on frequently accessed fields. This supports efficient catalogue
search, transaction lookup, overdue monitoring, and demand analysis.

% ============================================================
% LIBRARY INTELLIGENCE AGENT
% ============================================================

\section{Library Intelligence Agent}

The Library Intelligence Agent provides natural-language interaction with
the library catalogue.

Its responsibilities include:

\begin{itemize}

    \item interpreting natural-language resource requests,

    \item extracting relevant keywords and categories,

    \item searching library resources through controlled application APIs,

    \item checking resource availability,

    \item providing availability-grounded recommendations,

    \item answering common resource-related queries, and

    \item logging unfulfilled searches for subsequent demand analysis.

\end{itemize}

The agent does not receive unrestricted access to the database. Database
operations are exposed through controlled application APIs and tools.

A conventional keyword-based search remains available as a fallback if the
AI service is unavailable or if the extracted intent does not meet the
required confidence threshold.

% ============================================================
% RECORD ONBOARDING AGENT
% ============================================================

\section{Record Onboarding Agent}

The Record Onboarding Agent assists administrators with bulk student and
member record imports.

The onboarding workflow consists of:

\begin{enumerate}

    \item Receiving structured records.

    \item Validating required fields and data types.

    \item Checking roll numbers and unique identifiers.

    \item Performing duplicate detection.

    \item Flagging potential duplicate records for review.

    \item Inserting validated records through controlled backend
    operations.

    \item Generating an import summary.

\end{enumerate}

Potential duplicate records are flagged for review rather than being
automatically deleted.

% ============================================================
% LIBRARY MONITORING AGENT
% ============================================================

\section{Library Monitoring Agent}

The Library Monitoring Agent provides scheduled monitoring and
resource-management assistance.

Its responsibilities include:

\begin{itemize}

    \item detecting resources approaching their due date,

    \item detecting overdue resources,

    \item calculating applicable overdue fines,

    \item generating reminder and notification records,

    \item analysing issue frequency and resource availability, and

    \item identifying high-demand and low-availability resources for
    purchase consideration.

\end{itemize}

Deterministic operations such as due-date comparison and fine calculation
are performed by the backend and database rather than delegated to a
language model.

AI functionality is used where analysis, recommendation, or summarisation
provides additional value.

% ============================================================
% AI SAFETY
% ============================================================

\section{AI Safety and Control}

The AI layer is intentionally constrained.

The agents do not receive unrestricted SQL access. Instead, database
operations are exposed through predefined backend APIs and controlled
tools.

The following safeguards are incorporated:

\begin{itemize}

    \item parameterized database queries,

    \item role-based authorization,

    \item structured AI outputs,

    \item controlled tool invocation,

    \item deterministic handling of financial and circulation operations,

    \item backend validation before database modification, and

    \item conventional search fallback when AI services are unavailable.

\end{itemize}

% ============================================================
% CI/CD
% ============================================================

\section{CI/CD and Deployment}

The development process incorporates continuous integration and continuous
delivery practices.

The CI pipeline is intended to:

\begin{enumerate}

    \item install project dependencies,

    \item perform linting and static checks,

    \item execute automated tests,

    \item validate the application build, and

    \item deploy successful builds to the configured environment.

\end{enumerate}

This allows each development increment to be independently tested and
reduces the risk of integration failures during later development stages.

% ============================================================
% CHAPTER 3: PROTOTYPE IMPLEMENTATION AND EVALUATION
% ============================================================

\chapter{Prototype Implementation and Evaluation}

\section{Prototype Implementation}

The prototype establishes the major architectural components required by
the project proposal.

The implementation is organized into:

\begin{itemize}

    \item frontend web interface,

    \item backend REST API layer,

    \item MySQL database,

    \item controlled AI agent layer,

    \item automated monitoring workflow, and

    \item testing and deployment infrastructure.

\end{itemize}

The system is designed so that conventional library operations remain
functional independently of the AI layer.

\section{Core Library Workflow}

The primary student workflow consists of:

\begin{enumerate}

    \item Student authentication.

    \item Resource search.

    \item Availability verification.

    \item Resource issue.

    \item Due-date tracking.

    \item Resource return.

    \item Fine calculation where applicable.

\end{enumerate}

The administrator workflow additionally provides:

\begin{enumerate}

    \item Student/member management.

    \item Resource management.

    \item Bulk record onboarding.

    \item Overdue monitoring.

    \item Notification management.

    \item Demand analysis.

    \item Purchase recommendations.

\end{enumerate}

\section{Intelligent Resource Discovery Workflow}

The intelligent-search workflow follows the sequence:

\begin{enumerate}

    \item The student submits a natural-language query.

    \item The Library Intelligence Agent extracts the intent, keywords,
    and category.

    \item The backend performs controlled catalogue retrieval.

    \item Resource availability is calculated using current transaction
    data.

    \item Relevant results are ranked and returned.

    \item The search is logged for subsequent demand analysis where
    appropriate.

\end{enumerate}

The AI layer assists with interpretation and recommendation, while actual
availability information is retrieved from the database.

\section{Record Onboarding Workflow}

The bulk onboarding workflow is:

\begin{enumerate}

    \item The administrator uploads a structured batch of records.

    \item The Record Onboarding Agent validates the records.

    \item Duplicate and potentially duplicate records are identified.

    \item Valid records are passed to the backend.

    \item The backend executes insertion as a controlled transaction.

    \item The system returns an import summary containing accepted and
    skipped records.

\end{enumerate}

\section{Monitoring and Notification Workflow}

The monitoring workflow periodically evaluates active issue records.

For each active transaction:

\begin{enumerate}

    \item The system compares the current date with the due date.

    \item Overdue transactions are identified.

    \item Overdue days are calculated.

    \item Applicable fine records are created or updated.

    \item Notification records are inserted into the notification queue.

    \item Demand information is subsequently analysed using search and
    circulation records.

\end{enumerate}

% ============================================================
% TESTING
% ============================================================

\section{Testing Methodology}

The prototype follows a layered testing approach.

\subsection{Unit Testing}

Individual backend functions and AI-agent processing components are tested
independently using controlled inputs and expected outputs.

Particular attention is given to:

\begin{itemize}

    \item input validation,

    \item intent extraction,

    \item database query functions,

    \item availability calculation,

    \item fine calculation, and

    \item duplicate detection.

\end{itemize}

\subsection{Integration Testing}

Integration tests verify interactions between the frontend, backend APIs,
database, and controlled AI layer.

Representative workflows include:

\begin{itemize}

    \item natural-language search to catalogue retrieval,

    \item resource selection to issue transaction,

    \item bulk onboarding to database insertion, and

    \item overdue detection to notification generation.

\end{itemize}

\subsection{Performance Testing}

Performance testing measures API response times and resource-discovery
latency under controlled test conditions.

The database is evaluated for common catalogue and circulation queries,
with indexing and pagination used to support response-time requirements.

\subsection{Functional Testing}

Functional testing verifies that the principal student and administrator
workflows behave according to their expected requirements.

% ============================================================
% EVALUATION CRITERIA
% ============================================================

\section{Evaluation Criteria}

The prototype is evaluated using measurable workflow, accuracy,
performance, and reliability criteria.

\subsection{Primary Metrics}

\begin{table}[H]
\centering
\caption{Primary Prototype Evaluation Metrics}
\label{tab:primarymetrics}

\begin{tabularx}{\textwidth}{>{\raggedright\arraybackslash}p{3.5cm}
                              X
                              >{\centering\arraybackslash}p{2.5cm}}
\toprule
\textbf{Metric} & \textbf{Measurement} & \textbf{Target} \\
\midrule

Resource discovery time &
Time required to locate a requested or suitable available resource for
predefined queries &
$\leq 10$ seconds \\

Duplicate detection precision &
Proportion of records flagged as duplicates that are actual duplicates
in a labelled dataset &
$\geq 90\%$ \\

Overdue detection coverage &
Proportion of known overdue transactions correctly detected &
$\geq 99\%$ \\

Demand recommendation accuracy &
Proportion of high-demand and low-availability resources correctly
identified according to predefined thresholds &
$\geq 90\%$ \\

\bottomrule
\end{tabularx}
\end{table}

\subsection{Secondary Metrics}

\begin{table}[H]
\centering
\caption{Secondary Prototype Evaluation Metrics}
\label{tab:secondarymetrics}

\begin{tabularx}{\textwidth}{>{\raggedright\arraybackslash}p{3.5cm}
                              X
                              >{\centering\arraybackslash}p{2.5cm}}
\toprule
\textbf{Metric} & \textbf{Measurement} & \textbf{Target} \\
\midrule

API response time &
Median response time for common non-AI application requests &
$\leq 500$ ms \\

AI answer validity &
Proportion of predefined library questions answered consistently with
database information &
$\geq 90\%$ \\

Notification delay &
Time between an overdue condition becoming eligible and generation of
its notification &
Measured experimentally \\

System reliability &
Availability of the deployed application during the evaluation period &
$\geq 99\%$ \\

Test coverage &
Percentage of critical backend logic covered by automated tests &
$\geq 70\%$ \\

\bottomrule
\end{tabularx}
\end{table}

\section{Validation Strategy}

The evaluation uses controlled datasets and predefined scenarios.

For resource discovery, conventional search can be compared with
natural-language interaction by measuring the time required to locate
resources.

For onboarding, a labelled dataset containing valid, invalid, and
duplicate records can be used to evaluate duplicate-detection precision.

For monitoring, synthetic borrowing records with known due dates and
overdue states can be used to measure detection coverage and notification
latency.

For demand analysis, historical or synthetic issue records can be used to
define high-demand resources according to predetermined thresholds. The
resulting recommendations can then be compared against expected labels.

\section{Analysis}

The prototype is evaluated primarily on functional correctness,
integration, reliability, and measurable performance rather than on
novelty of the underlying AI models.

The architecture provides a clear separation between deterministic
operations and AI-assisted operations. This reduces the risk of
hallucinated or inconsistent information affecting critical library
transactions.

The principal engineering challenge is integrating AI-assisted features
without making the core system dependent on an external AI service.
Maintaining conventional search and deterministic backend workflows as
fallbacks addresses this concern.

The evaluation targets provide a basis for measuring whether the prototype
meets the intended software-engineering objectives. Actual achieved values
should be recorded after executing the corresponding controlled test
datasets and performance experiments.

% ============================================================
% CHAPTER 4: SCALABILITY
% ============================================================

\chapter{Scalability and Deployment}

\section{Scalability}

The system is designed to scale in both theoretical and practical terms.

\subsection{Logical Scalability}

\begin{itemize}

    \item Frontend and backend responsibilities are separated.

    \item Backend APIs are designed to be stateless where practical.

    \item Database queries use indexing and pagination for frequently
    accessed records.

    \item AI agents are modular and can be extended or replaced without
    rewriting the complete application.

\end{itemize}

\subsection{Operational Deployability}

\begin{itemize}

    \item The application can be containerized using Docker.

    \item A managed relational database can be used for deployment.

    \item CI/CD supports repeatable staging and production deployments.

    \item The architecture can support a larger number of users and
    records without requiring a fundamental redesign.

\end{itemize}

\section{Engine Availability}

The project relies on mature and readily available software engines rather
than requiring fundamental algorithmic research.

The required engineering components include:

\begin{itemize}

    \item a modern web framework for frontend and backend development,

    \item REST API infrastructure for application integration,

    \item MySQL or another managed relational database,

    \item authentication and role-based authorization libraries,

    \item mature LLM or AI APIs and supporting agent frameworks,

    \item automated testing frameworks,

    \item GitHub Actions or an equivalent CI runner, and

    \item Docker and a platform-hosted or cloud deployment target.

\end{itemize}

\section{Deployment Strategy}

The deployment strategy follows the project's incremental development
approach.

The application can first be deployed to a staging environment for
validation. Successful builds can subsequently be promoted to the
production environment through the CI/CD pipeline.

The separation between frontend, backend, database, and AI components
allows individual services to be scaled according to workload.

% ============================================================
% RISKS AND MITIGATIONS
% ============================================================

\section{Risks and Mitigations}

\begin{itemize}

    \item \textbf{AI hallucination or incorrect resource information:}
    AI responses are grounded in current database records through
    controlled retrieval tools and predefined APIs.

    \item \textbf{Excessive project scope:}
    The core library workflow is completed before AI features are added,
    and later functionality is delivered incrementally.

    \item \textbf{Duplicate detection errors:}
    Potential duplicates are flagged for human review instead of being
    automatically deleted.

    \item \textbf{AI service unavailability:}
    Conventional keyword-based resource search remains available as a
    fallback.

    \item \textbf{Unauthorized automated actions:}
    Agents operate through role-controlled APIs, and irreversible
    operations require explicit backend authorization.

    \item \textbf{Insufficient evaluation data:}
    Synthetic and labelled library datasets can be prepared so that
    evaluation does not depend entirely on live institutional data.

    \item \textbf{Deployment instability:}
    A staging environment and CI/CD pipeline make deployments repeatable.

    \item \textbf{Database performance degradation:}
    Frequently accessed fields are indexed and database operations use
    pagination and optimized queries.

\end{itemize}

% ============================================================
% CHAPTER 5: CONCLUSION
% ============================================================

\chapter{Conclusion}

\section{Summary of Work}

The prototype establishes the foundation of an AI-enabled Smart Library
Management System that combines conventional library workflows with
controlled AI-assisted functionality.

The system addresses four principal areas:

\begin{enumerate}

    \item resource discovery,

    \item student/member record onboarding,

    \item overdue and notification management, and

    \item resource demand analysis.

\end{enumerate}

The architecture separates the frontend, backend, database, and AI agent
layers, allowing individual components to be developed and tested
independently.

\section{Key Findings}

The prototype demonstrates how AI can be incorporated into a conventional
software system without allowing the AI layer to directly control critical
database operations.

The appropriate separation of responsibilities is:

\begin{itemize}

    \item \textbf{Deterministic backend/database operations:}
    availability calculation, issue and return transactions, due-date
    comparison, and fine calculation.

    \item \textbf{AI-assisted operations:}
    natural-language interpretation, resource recommendation,
    duplicate-detection assistance, demand analysis, and summarisation.

\end{itemize}

This separation improves reliability while still providing useful
AI-assisted interaction.

\section{Future Work and Recommendations}

Future development can extend the prototype through:

\begin{enumerate}

    \item improved semantic and vector-based resource search,

    \item more advanced recommendation ranking,

    \item improved duplicate-detection and entity-resolution methods,

    \item richer librarian analytics,

    \item improved demand forecasting,

    \item automated notification delivery,

    \item support for additional resource types beyond books,

    \item stronger authentication and role-management mechanisms, and

    \item deployment under a production cloud environment.

\end{enumerate}

\section{Final Remarks}

The project follows the software-engineering objective of delivering
practical value through incremental development rather than attempting to
solve an open-ended AI research problem.

The prototype provides a foundation for a deployable library-management
platform in which conventional software handles critical deterministic
operations and controlled AI agents provide natural-language interaction,
automation, and decision support.

The system can therefore be evaluated using measurable criteria covering
functionality, accuracy, performance, reliability, and scalability.

\end{document}
